\documentclass[11pt]{article}

\usepackage[margin=1in]{geometry}
\usepackage{amsmath,amssymb,amsthm}
\usepackage{enumitem}
\usepackage[colorlinks=true,linkcolor=blue,urlcolor=blue]{hyperref}

\newcommand{\MSO}{\ensuremath{\mathrm{MSO}}}
\newcommand{\MSOtwo}{\ensuremath{\mathrm{MSO}_2}}
\newcommand{\Bags}{\ensuremath{\mathrm{BAGS}}}
\newcommand{\adh}{\ensuremath{\mathrm{adh}}}
\newcommand{\cone}{\ensuremath{\mathrm{cone}}}

\title{Issues in \texttt{lecture\_note\_expanded.tex}\\
  Found While Preparing the Lean Formalization}
\author{}
\date{Review of the revised (2026-06) version; July 5, 2026}

\begin{document}
\maketitle

\section*{Scope and summary}

This report supplements \texttt{lecture\_note\_expanded\_issues.tex}, which
reviewed the draft dated June 19, 2026 and whose findings were subsequently
repaired.  The present pass reviewed the revised
\texttt{lecture\_note\_expanded.tex}---the version incorporating all of those
repairs---while preparing its Lean~4 formalization, checking both fidelity to
\texttt{lecture\_note.pdf} and mathematical correctness.  The mathematical
content of the revision was confirmed to be sound except for one remaining
gap, described and repaired below.  It is precisely the gap anticipated by
the closing remark of the earlier report: the defining-pair and
defining-tuple machinery depended on a fully proved encoding--decoding
correspondence, and one half of that correspondence was still missing.

Further issues found during the formalization itself can be appended to this
report in the same format.

\section{The translation-correctness setup identifies the decoding of an
encoding with the original graph}

\paragraph{Location.} The paragraph ``Correctness of the translation''
preceding the translation-correctness lemma; implicitly, the statements of
the recognition and atomic lemmas (\(\varphi_{\mathrm{vtx},i}\),
\(\varphi_{\mathrm{vtx}}\), \(\varphi_{\mathrm{set}}\),
\(\varphi_{\mathrm{adj}}\), \(\varphi_Q\), \(\varphi_{\mathrm{eq}}\),
\(\varphi_{\mathrm{cont}}\)).

The recognition and atomic lemmas are stated for a legal tree and \emph{its
decoding}: a tuple satisfies \(\varphi_{\mathrm{vtx}}\) iff it is the defining
tuple of a vertex \emph{of the decoded graph}, and so on.  The
translation-correctness lemma, however, speaks about a given triple
\((G,c,(T,\beta))\) and its encoding \((T,\rho)\).  Its preamble bridged the
two by asserting that ``by the legal-iff-encoding lemma its decoded graph is
\(G\)''.  That lemma proves two things: every encoding is legal, and every
legal tree is the encoding of its own decoding---that is, re-encoding is a
left inverse of decoding.  It does not prove that decoding the encoding of a
given triple returns that triple.  Indeed it cannot do so literally: the
decoded vertices are \(\sim\)-classes of slots \(v(t,i)\), not the vertices
of \(G\), so the best possible statement is a canonical isomorphism, and that
isomorphism---which must also preserve colors and occurrence sets, so that
defining pairs and defining tuples on the two sides agree---had never been
constructed.  This is the surviving trace of the original note's Lemma~9,
whose unproved ``moreover'' clause (any triple encoding to \((T,\rho)\)
\emph{is} the decoding) the expanded note had dropped without replacement.

The gap is load-bearing.  In the quantifier cases of the
translation-correctness induction, the forward direction needs the defining
tuple of a witness vertex of \(G\) to satisfy the guard
\(\varphi_{\mathrm{vtx}}\), while the backward direction needs every tuple
satisfying the guard to be the defining tuple of a vertex of \(G\)---not
merely of the decoded graph.

\paragraph{Repair.} Two repairs are possible.  (i) Prove the missing half of
the correspondence: decoding is a left inverse of encoding, up to a canonical
isomorphism.  (ii) Restate the recognition lemmas relative to a given encoded
triple, proving them directly from the defining-pair lemma on
\((G,c,(T,\beta))\)---the letters of an encoding record the bag and adhesion
data of the original decomposition verbatim---and bypassing the decoding
altogether; this removes the decoding from the critical path of Theorem~1
entirely and is the natural route for a formalization.  For the note, option
(i) was chosen: it keeps the recognition lemmas in their general legal-tree
form and completes the encoding--decoding correspondence, which is of
independent interest.

\paragraph{Status (2026-07).} Resolved in \texttt{lecture\_note\_expanded.tex}
by option (i), in three steps.

First, a new lemma \emph{Decoding inverts encoding} is placed immediately
after the legal-iff-encoding lemma.  For the encoding \((T,\rho)\) of a triple
\((G,c,(T,\beta))\), with decoding \((\bar G,\bar c,(T,\bar\beta))\), the slot
map
\[
  \Phi\bigl([v(t,i)]\bigr) := \text{the unique vertex of color \(i\) in
  \(\beta(t)\)}
\]
is shown to be a well-defined isomorphism of
\(\tau_{\mathcal P}\)-graphs \(\bar G\cong G\) which (i) preserves colors and
(ii) matches bags: \(\Phi(w)\in\beta(t)\) iff \(w\in\bar\beta(t)\), so that
occurrence sets satisfy \(\Bags(\Phi(w))=U_w\).  Well-definedness on the
generators of \(\sim\) follows from bag-injectivity; surjectivity from vertex
coverage; injectivity from the running-intersection axiom combined with the
path characterization of \(\sim\) through the subgraphs \(\Gamma_i\); and
preservation and reflection of adjacency and of the unary predicates from the
fact that the letters of an encoding are the color-renamed induced
substructures of \(G\), with reflection of adjacency using the edge-coverage
axiom.  Together with the legal-iff-encoding lemma, encoding and decoding are
now mutually inverse: re-encoding is a left inverse of decoding on legal
trees, and decoding inverts encoding up to the canonical isomorphism
\(\Phi\).

Second, a new corollary \emph{Recognition over an encoding} closes the
atomic-formula subsection: for the encoding of \((G,c,(T,\beta))\), the
recognition and atomic lemmas hold verbatim with the decoding replaced by
\((G,c,(T,\beta))\) itself, defining pairs and tuples being taken in the
original triple.  The proof observes that \(\Phi\) matches colors and
occurrence sets, so the defining pair of a decoded vertex \(w\) equals the
defining pair of \(\Phi(w)\) and defining tuples coincide coordinatewise;
bijectivity of \(\Phi\) makes the tuples arising on the two sides literally
the same, and the atomic cases transfer along the
\(\tau_{\mathcal P}\)-isomorphism.

Third, the ``Correctness of the translation'' paragraph now invokes this
corollary instead of asserting that the decoded graph is \(G\), and the
atomic, existential-vertex, and existential-set cases of the
translation-correctness proof cite the corollary's items in place of the
decoding-level lemmas.

The two insertions renumber the subsequent statements of the note; in
particular, the cited effective \MSO{}--tree-automaton correspondence,
formerly Theorem~26, is now Theorem~28.

\paragraph{Remark for the formalization.} Although the note adopts repair
(i), the Lean development should still prefer route (ii) for the main
theorem: state and prove the recognition lemmas relative to a given encoded
triple, so that the slot-quotient decoding and the isomorphism \(\Phi\) are
not on the critical path of Theorem~1.  The new lemma then corresponds to the
\emph{decode\(\circ\)encode} half of the planned encode/decode equivalence
(stage 5 of \texttt{docs/PLAN.md}) and can be formalized independently.

\end{document}
